\section{Diagnostic Prototype Details}
\label{app:diagnostic-details}

\subsection{Controlled factorial closure}

The controlled experiment is generated from existing scikit-learn digits data
by
\begin{verbatim}
python -u experiments/optimizer_calculus/interaction_curvature_experiment.py `
  --output-dir experiments/optimizer_calculus/interaction-curvature-results
\end{verbatim}
The complete configuration is
\path{experiments/optimizer_calculus/interaction_curvature_config.json}.
The paper-supplied mirror under \path{artifacts/interaction_curvature} contains
the final configuration snapshot, Boolean responses, pair-curvature checks,
continuous held-out checks, split robustness, allocations, Gaussian campaign
summaries, all 4500 confirmatory empirical observations, bootstrap intervals,
finite-response traces, and the consolidated summary.

The dataset is restricted to digits 3 and 8 and stratified into 249 training
and 108 test examples.  Standardization is fit on the training data only.  An
intercept gives 65 hidden coordinates, all ridge-regularized with
$\lambda=0.15$.  The three additive losses emphasize positive-class samples,
negative-class samples, and central-occlusion robustness with scale $0.75$.
Newton solves use analytic gradients and Hessians, Armijo damping, tolerance
$10^{-11}$, and at most 80 iterations.  Pair curvature uses order-24
Gauss--Legendre quadrature; eight additional splits use order 16.

The Gaussian pilot allocates 420 main and 80 held-out replications.  Exact
noncentral-$\chi^2$ power calibration multiplies both by the first integer
reaching 0.95 power, which is nine.  Confirmatory simulation uses 100,000
campaigns.  The empirical channel uses minibatches of size 1024, 200 pilot
observations per mask, a disjoint confirmatory sample of 3780 main and 720
held-out observations, and 20,000 stratified bootstrap draws.  The pilot is
used only for variance estimation and allocation.

Finite-response tracking runs fixed-step gradient descent for horizons
$0,1,2,4,8,16,32,64,128,256,512,1024$.  At every horizon the script verifies
that the actual pair-effect error is below the sum of the four optimized-value
gaps, and that this exact gap bound is below the analytic strong-convexity
bound.

The verification commands are
\begin{verbatim}
python -m pytest `
  experiments/optimizer_calculus/test_interaction_curvature_experiment.py -q

python experiments/optimizer_calculus/verify_interaction_curvature_results.py `
  newpapers/03-geometric-nongeometric-optimizer-calculus/artifacts/interaction_curvature
\end{verbatim}
They return \texttt{6 passed} and
\texttt{interaction-curvature result audit passed}, respectively.

\subsection{Quadratic benchmark}

The deterministic quadratic benchmark uses three problem families:
\begin{enumerate}[leftmargin=*,itemsep=2pt]
\item rotated strongly convex quadratics;
\item diagonal-plus-low-rank quadratics;
\item block-coupled quadratics.
\end{enumerate}
The script records final objective gap, gradient-call count, and convergence
under a gradient-norm tolerance.  The default repository command is
\begin{verbatim}
python experiments/toy/gng_optimizer_benchmark.py --output experiments/toy/gng-results
\end{verbatim}
and the larger diagnostic run used during development was
\begin{verbatim}
python experiments/toy/gng_optimizer_benchmark.py --output experiments/toy/gng-results `
  --dim 64 --instances 4 --memory 8 --max-iter 220 --grad-tol 1e-8
\end{verbatim}
The paper reports the later optimized diagnostic checks because they use
\textsc{GNG-FullMetricProbe} v2, which removes an unnecessary final gradient
validation call and therefore uses $d+1$ gradient calls (65 when $d=64$), as
recorded in \path{experiments/toy/gng-results-optimized-64x1/gng_optimizer_summary.json}.
The saved diagnostic summary also contains \textsc{HeavyBall-oracle}; its poor
performance on these high-condition-number instances should not be read as a
general statement about heavy-ball methods.  The \textsc{Adam-tuned} entry in
that summary is selected from a six-point learning-rate sweep, whose tuning
cost is separate from the recorded optimizer-call count.

\subsection{Prototype descriptions}

\paragraph{GNG-FullMetricProbe.}
The method evaluates $g_0=\nabla f(x_0)$ and then evaluates gradients at
$x_0+e_i$ for $i=1,\ldots,d$.  For a deterministic quadratic, these probes
recover the columns of $H$.  The update $x_+=x_0-H^{-1}g_0$ is the exact
Newton step.  The final gradient can be computed internally as
$g_0+H(x_+-x_0)$, so no final oracle call is needed.

\paragraph{GNG-KrylovMetric.}
The method performs sequential directional probes to obtain $Hp$ along
search directions.  It uses exact quadratic line search and conjugacy-style
updates.  Its interpretation in this paper is a growing Krylov subspace
metric, not a new claim about conjugate-gradient theory.
The reported implementation is not presented as a standard CG baseline.

\paragraph{RSDLR-BFGS-Hybrid.}
The method combines a bounded diagonal inverse metric with a limited secant
memory.  It keeps a mixture of recent secant pairs and pairs with high
residual under the current diagonal model.  In the present benchmark it did
not outperform FIFO-LBFGS.

\subsection{Formal trace audit}

The trace-level residual audit in \cref{sec:experiments} is produced
by the scripts under \texttt{experiments/lane\_audit}.  Each worker runs
\texttt{optimizer\_trace.py}, records the raw gradient before the optimizer
step, applies the optimizer, records the parameter update, and writes
\texttt{trace\_residual\_summary.json}, \texttt{trace\_residuals.csv}, and a
GER plot.  In the rich runs, each worker also writes bounded-diagonal
sensitivity metrics, alternative visible-covector metrics, layer-scalar
metrics, penalized layer-scalar path-fit curves, partial CSV files,
progress JSON, and explicit checkpoints.  The queued batch is launched by
\texttt{lane\_orchestrator.py}.  The paper tables and figures are generated
by:
\begin{verbatim}
python experiments/lane_audit/analyze_formal_run.py `
  --run-dir tmp/remote_runs/paper03_rich_family_b1024_l28_20260710_2223 `
  --output-dir newpapers/03-geometric-nongeometric-optimizer-calculus/figures `
  --prefix family_sensitivity

python experiments/lane_audit/analyze_formal_run.py `
  --run-dir tmp/remote_runs/paper03_rich_cifar10_b1024_l28_20260710_2223 `
  --output-dir newpapers/03-geometric-nongeometric-optimizer-calculus/figures `
  --prefix cifar10_audit

python experiments/lane_audit/analyze_paper03_rich_runs.py `
  --run-dir tmp/remote_runs/paper03_rich_family_b1024_l28_20260710_2223 `
            tmp/remote_runs/paper03_rich_cifar10_b1024_l28_20260710_2223 `
  --output-dir newpapers/03-geometric-nongeometric-optimizer-calculus/figures `
  --prefix paper03_rich
\end{verbatim}

Both formal batches use lane count $28$, batch size $1024$, checkpoint and
partial-write interval $256$, and $4096$ recorded optimizer steps per job.
The formal jobs did not download datasets on the remote hosts.  The bounded
diagonal family is $A_{ii}\in[10^{-6},1]$ unless the sensitivity grid names a
different lower or upper bound.  Raw TGER is the $\ell_2$ trace explanation
rate when each step chooses its own map.  Coherent TGER is the no-variation
$B_{\rm geo}=0$ specialization: one shared map is fit to the whole trace.  The
coherent residual is computed by the closed-form streaming accumulator in
\texttt{optimizer\_trace.py}; it solves the same least-squares problem as
stacking all updates and gradients while avoiding full trace-vector storage
in host memory.  In the audit script, the names \texttt{muon} and
\texttt{gng\_muon} both call the local \textsc{GNGMuon} implementation with
different hyperparameter settings.  They are retained as implementation
labels for the saved summaries, but they should not be interpreted as
official reference implementations of external Muon optimizers.

The family-sensitivity rich run is
\texttt{paper03\_rich\_family\_b1024\_l28\_20260710\_2223}.  It uses datasets
\texttt{mnist,fashion\_mnist}, optimizers
\texttt{adamw,sophia,muon,gng\_muon}, seeds \texttt{7,8,9}, width $32$,
train subset $8192$, and test subset $2048$.  It contains $24$ jobs,
$98304$ optimizer steps, and approximately $1.01\times10^8$ sample views.  It
ran on one NVIDIA RTX PRO 6000 Blackwell Server Edition GPU; the batch
monitor recorded $2990$ seconds of wall time, mean GPU utilization $99.5\%$,
peak GPU utilization $100\%$, and peak VRAM $33015$ MiB.

The CIFAR-10 rich run is
\texttt{paper03\_rich\_cifar10\_b1024\_l28\_20260710\_2223}.  It uses dataset
\texttt{cifar10}, the full $50000/10000$ train/test split, width $48$,
optimizers \texttt{adamw,sgd,lion,sophia,muon,gng\_muon}, and seeds
\texttt{7,8,9,11}.  It contains $24$ jobs, $98304$ optimizer steps, and
approximately $1.01\times10^8$ sample views.  It ran on one NVIDIA RTX PRO
6000 Blackwell Server Edition GPU; the batch monitor recorded $3554$ seconds
of wall time, mean GPU utilization $95.7\%$, peak GPU utilization $100\%$,
and peak VRAM $53373$ MiB.

\subsection{Follow-up trace audits}

The runtime-constrained follow-up tables in
\cref{tab:followup-sector-geometry,tab:followup-interventions,tab:followup-tinyresnet-geometry}
are generated by:
\begin{verbatim}
$SECTOR="tmp/remote_runs/<sector-followup-run>"
$TINY="tmp/remote_runs/<tinyresnet-followup-run>"
$OUT="newpapers/03-geometric-nongeometric-optimizer-calculus/figures"

python experiments/lane_audit/analyze_paper03_followup_runs.py `
  --sector-run-dir $SECTOR `
  --tiny-run-dir $TINY `
  --output-dir $OUT `
  --prefix paper03_followup
\end{verbatim}
The analyzer checks that the sector run has $30$ optimizer/seed rows, the
tiny-ResNet run has $18$ rows, the final short-horizon step counts are
$128$ and $96$, and all reported accuracy, bounded-diagonal, channel-scalar,
layer-scalar, coherent, and alternative-covector metrics are present.

\begin{table}[t]
\centering
\caption{Synthetic residual sanity check used before the neural-network
follow-up audits.  Raw TGER allows the explaining geometry to vary by step;
coherent TGER uses one shared explaining map across the trace.}
\label{tab:appendix-followup-synthetic-sanity}
\small
% Generated by experiments/lane_audit/analyze_paper03_followup_runs.py
\begin{tabular}{lrrrr}
\toprule
Mechanism & Mean GER & Raw TGER & Coh. TGER & Gap \\
\midrule
constant diagonal & 1.000 & 1.000 & 1.000 & 0.000 \\
varying diagonal & 1.000 & 1.000 & 0.520 & 0.480 \\
momentum memory & 0.484 & 0.455 & 0.227 & 0.227 \\
decoupled decay & 0.925 & 0.924 & 0.761 & 0.163 \\
operator mixing & 0.785 & 0.785 & 0.494 & 0.291 \\
\bottomrule
\end{tabular}

\end{table}

The sector/intervention run is:
\begin{center}
\small\path{paper03_followup_sector_cifar_b1024_l28_20260711_0242}
\end{center}
It uses CIFAR-10 with the full $50000/10000$ split, width $48$, the
\texttt{small\_conv} model, batch size $1024$, lane count $28$, seeds
$7,8,9$, and ten optimizer settings: AdamW, AdamW/no weight decay,
SGD/Nesterov, SGD/no momentum, Lion, Lion/no sign, Sophia,
Adafactor-style, Muon-style, and GNG-Muon.
Each job records $128$ optimizer steps with partial trace writes and
checkpoints every $32$ steps.  The batch contains $30$ jobs and $3840$
recorded optimizer steps.  Its monitor recorded $2569$ seconds of batch wall
time, mean GPU utilization $91.4\%$, peak GPU utilization $100\%$, and peak
VRAM $58757$ MiB.

The tiny-ResNet follow-up run is:
\begin{center}
\small\path{paper03_followup_tinyresnet_cifar_b1024_l28_20260711_0242}
\end{center}
It uses CIFAR-10 with the full $50000/10000$ split, width $24$, the
\texttt{tiny\_resnet} model, batch size $1024$, lane count $28$, seeds
$7,8,9$, and six optimizer settings: AdamW, SGD/Nesterov, Lion, Sophia,
Adafactor-style, and Muon-style.  Each job records $96$ optimizer steps with
partial trace writes and checkpoints every $32$ steps.  The batch contains
$18$ jobs and $1728$ recorded optimizer steps.  Its monitor recorded $2482$
seconds of batch wall time, mean GPU utilization $99.7\%$, peak GPU
utilization $100\%$, and peak VRAM $49499$ MiB.

The follow-up campaign initially tested longer horizons, but calibration with
the channel-scalar and coherent covector diagnostics projected runtimes beyond
the intended two-hour-per-task budget.  The final reported follow-up runs
therefore keep the full optimizer/seed/geometry/covector matrix and shorten
the trace horizon.  They are used as mechanism and model-extension audits,
not as replacement performance benchmarks for the longer rich trace runs.
